\documentclass[a4paper]{article}
\usepackage[pdftex]{graphicx}
\usepackage[utf8]{inputenc}
\usepackage{enumerate}
\usepackage{siunitx}
\sisetup{locale=DE}
\usepackage{href-ul}
\hypersetup{
	colorlinks=true,
	linkcolor=blue,
	urlcolor=blue}
\usepackage{icomma}
\usepackage{amssymb}
\usepackage{geometry}
\geometry{a4paper, top=15mm, left=15mm, right=15mm, bottom=15mm,
	headsep=10mm, footskip=12mm}

\begin{document}
	
	\begin{enumerate}[1.]
		
		\item {\bf oscillations; thread pendulum}\\
		In the case of a pendulum, the connection between pendulum length and frequency was measured:
		
		\begin{tabular}{cccccc}
			$l$ in m & 0.20 & 0.40 & 0.60 & 0.80 & 1.0 \\
			$f$ in Hz & 1.1 & 0.79 & 0.64 & 0.56 & 0.50
		\end{tabular}
		
		\begin{enumerate}[a)]
			\item Draw a diagram from which the connection between the length $l$ and the \underline{oscillation period} $T$ can be seen.\\
			(x-axis= $ ~ l $-axis; $\SI{10}{\centi\meter}~\hat{=}~\SI{1}{\meter}$ ;
			y-axis=$ ~T$-axis: $\SI{5}{\centi\meter}~\hat{=}~\SI{1}{\second}$)
			\item How long is a pendulum that has an oscillation period of $\SI{1.5}{\second}$? Use your diagram.
			\item On what properties of the pendulum does the period of oscillation depend?
			\item A pendulum has a frequency of $f=\SI{0.5}{\hertz}$. How many periods pass in $\SI{24}{\second}$?
		\end{enumerate}
		
		
		{\bf \item The oscillating mass of a spring pendulum is $\SI{440}{\gram}$.} The oscillation occurs with a maximum speed of the pendulum body of $\SI{0.30}{\meter\over\second}$ . In the experiment, $\SI{12}{\second}$ is required for 10 oscillations. The spring hardness is $\SI{12}{\newton\over\meter}$. At time $t = \SI{0}{\second}$ the spring pendulum swings in the downward direction through the rest position
		
		\begin{enumerate}[a)]
			\item When is a vibration called a harmonic vibration? Give 2 characterizations for it.
			\item The oscillation period of the system can be calculated from the information in two ways. Compare the results and state the frequency.
			\item Set up the function $v(t)$ for the speed of the pendulum bob and calculate it at time $ t=\SI{7,9}{\second}$. \\
			Sketch the course of the graph of $v(t)$ in the interval $[0;\SI{1,2}{\second}]$. Use colored arrows to mark the times of the greatest acceleration amounts.
			\item Explain the approach we used in class to establish a connection between the amplitude and the maximum speed and use it to calculate the amplitude of the oscillation.\\
			$[\textnormal{False replacement result}: \hat{x}=\SI{5,0}{\centi\meter}]$
			\item Specify the magnitude and direction of the restoring force at the upper reversal point.
			
			\item By replacing the spring with the same mass, the oscillation period increases to $T=\SI{1.5}{\second}$. \\
			What spring hardness does the new spring have?
		\end{enumerate}
		
		{\bf \item Specify a change in a spring pendulum that causes...}
		
		\begin{enumerate}[a)]
			\item ...the oscillation period doubles
			\item ...the repulsive force is only half as great with the same deflection
		\end{enumerate}
		
	\end{enumerate}
	
	\fbox{
		\begin{minipage}{0.5\textwidth}
			For the solution, please \href{https://www.okuyakl.de/phys/p10penL407/le407.pdf}{click here} or scan the QR code.\\
			Additional worksheets are available at
			
			\href{http://www.okuyakl.com}{www.okuyakl.com}
		\end{minipage}
		\hfill
		\begin{minipage}{0.4\textwidth}
			\includegraphics[width=1.5 cm]{../../viecher/zwe03}
			\includegraphics[width=3 cm]{qre407}
			\includegraphics[width=2 cm]{../../viecher/afanticon1}
			
	\end{minipage}}
	
\end{document}